% Replacement abstract and introduction for the ISAC-assisted underserved-region
% detection and UAV-repositioning paper.
%
% Do not restore the previous claims about multi-area terrestrial-sensor data
% collection, hybrid mission time, edge-processing delay, or traditional-RAN
% versus O-RAN mission-time reduction: those are not evaluated by the current
% mathematical model or simulation workflow.
%
% Before final submission, add one concise quantitative result sentence to the
% abstract using only completed multi-seed experiments. Do not report the RF
% controller as evaluated until its grid reporter, training/inference pipeline,
% and probability-weighted K-means implementation are complete.
%
% The old multi-area TSD-collection figure does not match this paper and should
% be removed or redrawn. A replacement architecture figure should show ground
% UEs, TN gNBs, gNB-mounted UAV sensing nodes, E2/near-RT RIC fusion and service
% classification, the RF predictor, K-means centroids, and UAV movement. A
% suitable caption is: ``ISAC-assisted O-RAN closed loop for persistent
% repositioning-demand detection and gNB-mounted UAV control.''

\begin{abstract}
Low-altitude next-generation NodeB (gNB)-mounted Unmanned Aerial Vehicles
(UAVs) can provide flexible coverage when terrestrial infrastructure cannot
adequately serve spatially concentrated user equipment (UEs). However,
effective repositioning of UAVs requires the network to distinguish a genuine
coverage deficit from a transient attachment or handover condition and to
estimate where the underserved UEs are located without relying on simulator
ground-truth coordinates. This paper proposes an Integrated Sensing and
Communications (ISAC)-assisted Open Radio Access Network (O-RAN) framework for
detecting persistent UAV-repositioning demand. Each gNB-mounted UAV acts as a
system-level monostatic sensing node that models UE detection as a function of
echo Signal-to-Noise Ratio, produces noisy position observations, and supports
inverse-variance fusion at the near-real-time RAN Intelligent Controller. A
UAV is repositioned only when weak service persists and neither a feasible
terrestrial nor an already deployed UAV cell can resolve it. To support
proactive control, a Random Forest (RF) is formulated to predict future
grid-level repositioning demand, and its predicted probabilities are used in
probability-weighted K-means clustering. Across independent simulation seeds,
the proposed predictive controller reduces [METRIC] by [X]\% relative to
ISAC-reactive K-means while improving the fifth-percentile throughput by
[Y]\%. The resulting methodology separates the communication decision of
which UEs require intervention from the ISAC estimate of where that demand is
located.
\end{abstract}

\section{Introduction}
\label{sec:introduction}

Integrated Sensing and Communications (ISAC) is vital to future wireless
systems because it enables radio infrastructure to obtain environmental
information while providing communication service. Co-locating the two
functions can reduce
hardware duplication and improve the use of spectrum and energy compared
with independent sensing and communication platforms
\cite{ahmed2025advancements}. Unmanned Aerial Vehicles (UAVs) are particularly
attractive ISAC platforms because their three-dimensional mobility permits
sensing and coverage resources to follow changing traffic distributions and
temporarily underserved regions \cite{meng2023uav,mu2023uav}.

The benefit of this mobility depends on deciding both \emph{when} movement is
needed and \emph{where} the UAV should move. A low serving-cell Reference
Signal Received Power (RSRP) does not necessarily identify a genuine coverage
deficit. The UE may still be attaching, undergoing handover, or able to obtain
adequate service from another terrestrial or already deployed UAV cell.
Immediately repositioning a UAV in these cases can create unnecessary travel
and destabilise coverage elsewhere. At the same time, supplying perfect UE
coordinates to a positioning algorithm bypasses missed detections,
localisation uncertainty, and observation fusion, thereby overstating the
capability of ISAC-assisted control. These considerations require a clear
separation: communication measurements determine \emph{which} UEs represent
unresolved service demand, while ISAC observations estimate \emph{where} that
demand is located.

The Open Radio Access Network (O-RAN) architecture provides a programmable
control framework for combining these information sources. Through E2
reporting and xApp logic, the near-real-time RAN Intelligent Controller
(near-RT RIC) can process radio measurements, fused sensing observations,
cell state, and UAV-mobility decisions in a closed loop. The terrestrial
network is considered first, existing UAV coverage provides a fallback, and
new UAV movement is requested only when the available cells cannot resolve
persistent weak service. In this study, O-RAN and E2 operation are represented
inside ns-3 using the ns-O-RAN simulation model. An external production
near-RT RIC and a network-facing E2 Application Protocol deployment are
outside the present scope.

% IMPORTANT: replace the old TSD/multi-area artwork with a diagram that actually
% shows the control loop described by this caption and the following paragraph.
\begin{figure}[t]
    \centering
    \includegraphics[width=\linewidth]{Figures/fig1.png}
    \caption{ISAC-assisted O-RAN control loop for detecting persistent
    repositioning demand and controlling gNB-mounted UAVs.}
    \label{fig1}
\end{figure}

As illustrated in Fig.~\ref{fig1}, the considered system comprises ground
UEs, terrestrial-network (TN) gNBs, and gNB-mounted UAVs. Each gNB-mounted UAV
also carries a co-located transmit--receive ISAC subsystem. At every
controller epoch, sensing is attempted for all relevant UEs before
communication-based service classification. A system-level monostatic radar
equation determines echo Signal-to-Noise Ratio (SNR), a logistic model samples
detection, and an SNR-dependent Gaussian model generates noisy horizontal
position observations. When several UAVs detect the same UE, the near-RT RIC
computes an inverse-variance fused position. Simulator ground-truth
coordinates are restricted to observation generation, offline supervised
labels, and evaluation; they are not exposed to the online ISAC controller.

Following sensing, the controller examines the serving RSRP and the best
fresh, capacity-feasible TN and existing-UAV candidates. Runtime handover
selection retains TN priority. Missing attachment or a stale serving
measurement places a UE in an UNKNOWN state instead of immediately declaring
it underserved. A weak or timed-out UNKNOWN state with no feasible
alternative must persist beyond the handover time-to-trigger and controller
response interval before the UE becomes a movement target. Reactive K-means
then clusters only these persistent targets. Its position input is the true UE
coordinate for the oracle baseline and the fused observation for the
ISAC-reactive method. Thus, weak service and UAV-repositioning demand remain
distinct quantities.

Reactive clustering responds to current demand and may therefore lag changes
in UE location or radio service. To support proactive repositioning, we
formulate a Random Forest (RF) predictor over spatial grid cells. Its features
describe sensed UE density, persistent repositioning demand and its share,
serving-RSRP and sensing-uncertainty aggregates, temporal demand change, and
distance to the nearest UAV. Future labels are generated offline using true
future positions and the same persistent no-feasible-alternative criterion
under one frozen RF-disabled reference policy. This avoids generating labels
under a proactive controller whose own UAV movements alter future RSRP and
therefore the prediction target. During online operation, the RF estimates a
future-demand probability for each cell, and these probabilities weight the
K-means objective. The RF is trained offline; K-means has no training stage
and is re-solved at every controller epoch.

The framework is realised in an integrated 5G-LENA and ns-O-RAN simulation
environment. The evaluation separates the effect of movement, localisation,
and prediction through four methods: a static-UAV baseline, an
oracle-position reactive K-means controller, an ISAC-sensed reactive K-means
controller, and the proposed ISAC--RF predictive controller. Performance is
assessed using sensing detection and localisation accuracy,
repositioning-target error, persistent service-state statistics,
fifth-percentile throughput, packet delivery ratio, and outage time.

The principal contributions of this work are as follows:
\begin{itemize}
    \item We develop a system-level ISAC sensing pipeline for gNB-mounted UAVs,
    comprising echo-SNR calculation, probabilistic detection, SNR-dependent
    localisation error, and multi-UAV inverse-variance fusion.

    \item We introduce a sense-before-classify O-RAN control procedure that
    separates experienced weak service from UAV-repositioning demand through
    fresh-RSRP and spare-capacity feasibility, TN-first alternative-cell
    checking, explicit UNKNOWN handling, and persistent no-alternative logic.

    \item We formulate policy-conditioned RF labels for future grid-level
    repositioning demand and use predicted cell probabilities in a weighted
    K-means objective, while preventing true UE coordinates from entering
    online control.

    \item We extend an integrated 5G-LENA and ns-O-RAN simulation environment
    with reproducible sensing, service-decision, centroid-assignment, mobility,
    handover, RSRP, and flow traces, and define static-UAV, oracle-position
    reactive, and ISAC-sensed reactive baselines for controlled comparison
    \cite{5glenaNR,garey2023ran}.
\end{itemize}

The remainder of this paper is organised as follows. Section~\ref{sec:RW}
reviews related work on UAV-assisted ISAC, O-RAN control, and learning-based
UAV positioning. Section~\ref{sec:UAV-ORAN} presents the system architecture
and control loop. Section~\ref{sec:mathematical-model} develops the sensing,
service-state, prediction, clustering, and movement models.
Section~\ref{closed_loop_model} describes the ns-3 implementation and
experimental configuration. Section~\ref{sec:diss} reports and discusses the
results, and Section~\ref{sec:conc} concludes the paper.
